\scanpage{644}
% 644-01
$\unsure{v|\sigma}$ dans $\mathrm{GL}(2,\widehat{\mathbb Z})/\operatorname{Im}\widehat{\Pi}_0^{\unsure{\circ}}$, \unsure{factorisé} dans $\mathcal G/\Pi_0$, est égale à celle de $u$. Donc l'assertion que $M^\circ(-\bar\jmath)$ est dans le noyau de $M\longrightarrow\mathcal G/\Pi_0$, équivaut \deleted{au fait}\inserted{à celle} que $M^\circ(0)$ est dans ce noyau, donc que l'on a un diagramme comm.
% 644-02
\diagram{\begin{tikzcd}[ampersand replacement=\&]
M(0) \arrow[r] \arrow[d,"{(\varepsilon_2,\varepsilon_3)}"'] \& M \arrow[d,"{(45)}"] \\
\mu\times\mu \arrow[r,"?"] \& \mathcal G/\Pi_0
\end{tikzcd}}
% 644-03
\margin{\deleted{\illegible{2}} NB \deleted{\unsure{commut.}; et} la \unsure{correction|condition} $\varepsilon_3$ n'intervient pas dans la déf\unsure{inition} (47) de $M\longrightarrow\mathcal G/\Pi_0$, cf plus loin.}
% 644-04
Notons que le composé
% 644-05
\diagram{\begin{tikzcd}[ampersand replacement=\&, column sep=large]
M \arrow[r] \arrow[dr] \& \mathcal G/\Pi_0 \arrow[r] \& (\deleted{\unsure{\mathrm{GL}}}\mathbb Z/4\mathbb Z)^*\simeq\{\pm1\} \\
\& \mathrm{GL}(2,\mathbb Z/4\mathbb Z)^{\unsure{*}} \arrow[ur] \& \deleted{\unsure{\mathrm{GL}(2,\mathbb Z),\mathbb Z}} \arrow[l,shift left=1ex] \arrow[u,shift left=1ex]
\end{tikzcd}}
% 644-06
n'est autre que $\varepsilon_2$. Notons aussi que par la compatibilité
% 644-07
\diagram{\begin{tikzcd}[ampersand replacement=\&]
\mathcal G \arrow[r,hook] \arrow[dr,"{\mathrm{can.}}"'] \& M \arrow[d,"{(45)}"] \\
\& \mathcal G/\Pi_0
\end{tikzcd}}
% 644-08
l'opération de $M$ sur $\mathcal G/\Pi_0$ déduite de l'hom (45) n'est autre que l'opération \deleted{\unsure{régulière}}\inserted{\unsure{de groupe top.|de quotients}}, des \unsure{automs} sur $\mathcal G$ en passant au quotient par $\Pi_0$, opération qui, pour $u\in M^!$ et en particulier sur $M(\unsure{\omega|0})\subset M^!$, est donnée par $\varepsilon_2(u)$ (cf p. \unsure{554|594}). Ainsi
% 644-09
\mspage{643}

\scanpage{645}
% 645-01
\mspage{625}
% 645-02
On voit que $M_{\deleted{\unsure{0|a}}}^{\prime\prime!}$ s'envoie dans le centre de $\unsure{\mathcal G/\Pi_0|\mathcal G[\Pi_0]}$, qui est contenu dans $\mathcal G^+/\Pi_0$, et en fait égal à \deleted{centre} l'image $\mu$ de $\Pi/\Pi_0$ dans $\mathcal G^+/\Pi_0$ (puisque le centre de $\mathcal G^+/\Pi\simeq\unsure{\mathfrak S_3|\mathcal G_3}$ est trivial). Donc on a une factorisation
% 645-03
\diagram{\begin{tikzcd}[ampersand replacement=\&]
M_{\deleted{\unsure{0|a}}}^{\prime\prime!} \arrow[r] \& \mu \arrow[r,hook,"{\mathrm{centre}}"'] \& \mathcal G/\Pi_0
\end{tikzcd}}
% 645-04
et il faudrait vérifier que cet hom $M^{\prime\prime!}\longrightarrow\mu$ n'est autre que $\varepsilon_3$, qui est donc trivial sur $M^{\circ!}$ et a fortiori sur $M^{\circ!}(0)$.
% 645-05
\inserted{En fait, on \unsure{reconnaît} que le noyau n'est pas $\operatorname{Ker}(\varepsilon_3\mid M^{\prime\prime!})$ \deleted{\unsure{précisé}} (qui contient en \unsure{fait|effet} $\mu$\unsure{;|!}) mais $M_0^{\prime\prime!}$ --- on va y \uline{revenir} plus bas.}
% 645-06
\deleted{Admettons provisoirement ce fait, \unsure{trouve|trouvons} donc une factorisation de
% 645-07
\diagram{\begin{tikzcd}[ampersand replacement=\&]
N_\rho^* \arrow[r] \& \deleted{\unsure{M}}\mathcal G/\Pi_0
\end{tikzcd}}
en
% 645-08
\diagram{\begin{tikzcd}[ampersand replacement=\&]
(*)\quad N_\rho^*/M^\circ(-\bar\jmath) \arrow[r] \& \mathcal G/\Pi_0
\end{tikzcd}}
et comme l'hom composé
% 645-09
\diagram{\begin{tikzcd}[ampersand replacement=\&]
\mathbb D_\rho \arrow[r,hook] \& N_\rho^*/M^\circ(-\unsure{\bar\jmath|\jmath}) \arrow[r] \& \mathcal G/\Pi_0
\end{tikzcd}}
% 645-10
est injectif, on voit que le noyau de (*) est un sous-groupe de $N_\rho^*/M^\circ(-\bar\jmath)$ dont l'intersection avec $\mathbb D_\rho$ est $\{1\}$ --- il est donc $1$, ou d'ordre $2$ (puisque}
% 645-11
\mspage{644}

\scanpage{646}
% 646-01
\deleted{$\mathbb D_\rho$ est d'indice $2$ dans $N_\rho^*/M^\circ(-\bar\jmath)$).
% 646-02
Je vois mieux --- \unsure{essayons} la
% 646-03
\uline{Proposition} Considérons la filtration de $\mathcal G^{\unsure{\wedge|*}}/\Pi_0^{\unsure{\wedge}}$ par $\mathcal G/\mathcal G^+\simeq\mu$, $\mathcal G^+/\Pi\simeq\unsure{\mathfrak S_3|\mathcal G_3}$, et $\Pi/\Pi_0\simeq\mu$. Considérons la filtration de $M$ par $M/M^{\prime\prime}\xrightarrow[\sim]{\varepsilon_2}\mu$, par $M^{\prime\prime}/M^{\prime\prime!}\simeq\unsure{\mathfrak S_3|\mathcal G_3}$, et $M^{\prime\prime!}/M^{\circ!}\overset{\varepsilon_3}{\simeq}\mu$.
% 646-04
Alors l'hom can.
\diagram{\begin{tikzcd}[ampersand replacement=\&]
M \arrow[r] \& \mathcal G/\Pi_0
\end{tikzcd}}
% 646-05
est compatible avec les filtrations, et induit \unsure{des} isomorphismes évidents sur les facteurs. En particulier le noyau de $M\longrightarrow\mathcal G/\Pi_0$ est $M^{\circ!}$.}
% 646-06
Admettant ce point, on trouve donc
% 646-07
\[
\operatorname{Ker}(N_\rho^*\longrightarrow\mathcal G/\Pi_0)
\unsure{\simeq|=} N_\rho^*\cap M_{\inserted{0}}^{\deleted{\unsure{\circ}}\inserted{\prime\prime}!}\deleted{\unsure{a}}
\]
% 646-08
et comme
\[
N_\rho^*\cap M_0^! = Z_{\rho0}^! = M^!(-\bar\jmath),
\]
% 646-09
on trouve que le noyau cherché est égal à :
\[
M^{\unsure{\prime|!}}(-\bar\jmath)\cap M^{\prime\prime}=M^\circ(-\bar\jmath),\quad\mathrm{cqfd}.
\]
% 646-10
Donc \deleted{l'hom} on a bien un hom injectif \unsure{(46)|(45)}, et comme les groupes
% 646-11
\mspage{645}

\scanpage{647}
% 647-01
\mspage{626}
% 647-02
en présence ont même ordre $24$, cet hom est bien un iso, cqfd.
% 647-03
d) Application à la structure de $\mathcal M_{\rho,\sigma}$ et à la définition d'un hom \deleted{$\unsure{\mathcal M}_{\rho,\sigma}$} $M\longrightarrow\mathcal M_{\rho,\sigma}/\mu_\rho$.
% 647-04
Rappelons que
\[
\tag{48}
\mathcal M_{\rho,\sigma}\overset{\mathrm{d\acute ef}}{=}
N_\rho^*\times_{\gamma_{\rho,\sigma}}N_\sigma^*\times_{\gamma_{\rho,\sigma}}G_{\rho,\sigma}
\]
% 647-05
où \deleted{$\unsure{\mathcal G}_{\rho,\sigma}$}
\[
\tag{49}
\left\{\begin{aligned}
G_{\rho,\sigma}&\simeq M=\underbrace{M(0)}_{\inserted{\simeq M_{0,\unsure{3|1}}^{\sim}(0)}}\cdot\mathcal G^{+\wedge}\\
\gamma_{\rho,\sigma}&\overset{\mathrm{d\acute ef}}{=}G_{\rho,\sigma}/\underbrace{\mathcal G^{+\wedge}}_{\inserted{\mathcal G_{\rho,\sigma}}}\simeq\unsure{\widetilde{\Gamma}_{\mathbb Q}|\widetilde{\Pi}_{\mathbb Q}}\ (\xleftarrow{\sim}M(0)),\\
N_\rho^*&=\operatorname{Norm}_M(L_\rho),\qquad N_\sigma^{\unsure{*}}=\operatorname{Norm}_M(L_\sigma)
\end{aligned}\right.
\]
% 647-06
D'autre part, nous avons introduit les sous-groupes remarquables
\[
\tag{50}
\left\{\begin{aligned}
Z_{\sigma0}^!&=M(\tfrac12)\subset N_\sigma^*\qquad(\mathrm{p.}\ 615\ (9))\\
\deleted{\unsure{A}}\inserted{N}_\rho^?&=\deleted{\unsure{\rho}}Z_\rho^!\cdot\mu_{\tau'}\subset N_\sigma^*\qquad(\mathrm{p.}\ 616\ (21)).
\end{aligned}\right.
\]
% 647-07
Le premier de ces groupes ne contient pas $\mu$, le deuxième le contient
\[
\tag{51}
\left\{\begin{aligned}
Z_{\sigma0}^!\cap\mu&=\{1\}\\
N_\rho^?\deleted{\unsure{0}}&\supset\mu
\end{aligned}\right.
\]
% 647-08
et
\diagram{(52)\quad\left\{\begin{array}{l}
\begin{tikzcd}[ampersand replacement=\&]
Z_{\sigma0}^! \arrow[r,"\sim"] \& \gamma_{\rho,\sigma}=\unsure{\widetilde{\Gamma}_{\mathbb Q}|\widetilde{\Pi}_{\mathbb Q}}
\end{tikzcd}\\[1.2em]
\deleted{\unsure{\operatorname{Ker}}\,N_\rho^?\unsure{=\mu|\longrightarrow\mu}}\quad
\begin{tikzcd}[ampersand replacement=\&, column sep=small]
{} \& {} \& {} \& \inserted{\unsure{\widetilde{\Gamma}_{\mathbb Q}|\widetilde{\Pi}_{\mathbb Q}}}\arrow[d,equal] \& \text{exacte}\arrow[d,bend left=25]\\
1\arrow[r] \& \mu\arrow[r] \& N_\rho^?\arrow[r] \& \gamma_{\rho,\sigma}\arrow[r] \& 1
\end{tikzcd}
\end{array}\right.}
% 647-09
\mspage{646}

\scanpage{648}

Soit alors
\[
\begin{array}{rcl}
(53)\qquad \mathcal{M}_{\rho,\sigma}^{\unsure{\natural|\sharp}}
&\subset&\mathcal{M}_{\rho,\sigma}\\[-2pt]
\Vert\,\mathrm{d\acute{e}f}&&\\[-2pt]
\unsure{N_{\rho}^{?}}\times_{\gamma_{\rho,\sigma}}
Z_{\sigma0}^{!}\times_{\gamma_{\rho,\sigma}}
\underbrace{G_{\rho,\sigma}}_{M}&&
\end{array}
\]

On a un \deleted{\unsure{hom.}}\inserted{plongement central} canonique
\[
\deleted{\epsilon\mapsto(1,\epsilon,1)}
\]
\diagram{
\begin{tikzcd}[ampersand replacement=\&, column sep=large]
(54)\qquad \mu\times\mu\times\mu \arrow[r,hook] \& \mathcal{M}_{\rho,\sigma}\\
(\epsilon,\epsilon',\epsilon'') \arrow[r,mapsto] \& (\epsilon,\epsilon',\epsilon'')
\end{tikzcd}}
provenant des plongements centraux
\deleted{\unsure{les}}
\diagram{
\begin{tikzcd}[ampersand replacement=\&, column sep=small]
\mu \arrow[r,hook] \& N_{\rho}^{*}, \& \mu \arrow[r,hook] \& N_{\sigma}^{*}, \&
\mu \arrow[r,hook] \arrow[dr,hook,bend right=25] \& G_{\rho,\sigma}\\
\&\&\&\&\& \mathcal{G}_{\rho,\sigma}^{+}\equiv\mathcal{G}^{+}\unsure{)} \arrow[u,hook,bend right=12]
\end{tikzcd}}
et on a
\[
\mathcal{M}_{\rho,\sigma}^{\unsure{\natural|\sharp}}
\cap(\mu\times\mu\times\mu)=(1\times\mu\times\mu)
\]
donc
\diagram{
\begin{tikzcd}[ampersand replacement=\&]
(55)\qquad \mu\times\mu \arrow[r,hook] \& \mathcal{M}_{\rho,\sigma}^{\unsure{\natural|\sharp}}.
\end{tikzcd}}

On désigne par $\mu_{\rho}$ le sous-groupe central $\mu$ de $N_{\rho}$ ou de $N_{\rho}^{?}$,
ou $\mu\times1$ dans $\mathcal{M}_{\rho,\sigma}^{\unsure{\natural|\sharp}}$.
Il est clair alors par (52), que l'on a une suite exacte
\diagram{
\begin{tikzcd}[ampersand replacement=\&, column sep=large]
(55)\qquad 1 \arrow[r] \& \mu_{\rho} \arrow[r,hook] \&
\mathcal{M}_{\rho,\sigma}^{\unsure{\natural|\sharp}} \arrow[r] \&
G_{\rho,\sigma} \arrow[r] \arrow[d,phantom,"\simeq"{rotate=90}] \& 1\\
\&\&\&M\&
\end{tikzcd}}
d'où des \unsure{hom.} canoniques

\mspage{647}

\scanpage{649}
\mspage{627}

\diagram{
\begin{tikzcd}[ampersand replacement=\&, column sep=large]
(56)\qquad M \arrow[r,"\sim"] \&
\mathcal{M}_{\rho,\sigma}^{\unsure{\natural|\sharp}}/\mu_{\rho} \arrow[r,hook] \&
\mathcal{M}_{\rho,\sigma}/\mu_{\rho}.
\end{tikzcd}}

\margin{On pose donc
$\Gamma_{\rho,\sigma}^{\unsure{\natural|\sharp}}
\overset{\mathrm{d\acute{e}f}}{=}
\mathcal{M}_{\rho,\sigma}^{\unsure{\natural|\sharp}}/\mathcal{G}_{\rho,\sigma}^{+}
\subset\Gamma_{\rho,\sigma}=\mathcal{M}_{\rho,\sigma}/\mathcal{G}_{\rho,\sigma}^{+}$.}

\unsure{Évidemment} le sous-groupe
$\mathcal{M}_{\rho,\sigma}^{\unsure{\natural|\sharp}}$ de $\mathcal{M}_{\rho,\sigma}$
contient $\mathcal{G}_{\rho,\sigma}=SG_{\rho,\sigma}\unsure{\simeq}\mathcal{G}^{+\wedge}$,
en fait on a un diagramme de suites exactes
\diagram{
\begin{tikzcd}[ampersand replacement=\&, column sep=small,row sep=large]
1 \arrow[r] \& \mu_{\rho}\times\mathcal{G}_{\rho,\sigma}^{\deleted{\unsure{+}}}
\arrow[r] \arrow[d,hook] \&
\mathcal{M}_{\rho,\sigma}^{\unsure{\natural|\sharp}} \arrow[r] \arrow[d] \&
\gamma_{\rho,\sigma} \arrow[r] \arrow[d,equal] \&1\\
(57)\quad1 \arrow[r] \&
\underbrace{SZ_{\rho}}_{L_{\rho}}\times\underbrace{SZ_{\sigma}}_{L_{\sigma}}
\times\mathcal{G}_{\rho,\sigma}^{+} \arrow[r] \&
\mathcal{M}_{\rho,\sigma} \arrow[r] \& \unsure{\gamma_{\rho,\sigma}|\Gamma_{\rho,\sigma}} \arrow[r] \&1
\end{tikzcd}}

Dans l'isomorphisme (56), le sous-groupe $\mathcal{G}^{+\wedge}$ de $M$
correspond au sous-groupe
$\mu_{\rho}\times\mathcal{G}_{\rho,\sigma}^{+}/\mu_{\unsure{\sigma|\rho}}
\simeq\mathcal{G}_{\rho,\unsure{\sigma}}$,
il s'envoie isomorphiquement sur le sous-groupe $\mathcal{G}_{\rho,\sigma}$
de $\mathcal{M}_{\rho,\sigma}/\mu_{\rho}$.

\deleted{Le sous-groupe $M(0)$ s'envoie}

On trouve ainsi
\diagram{
\begin{tikzcd}[ampersand replacement=\&, column sep=small,row sep=large]
\& \inserted{\unsure{\mathcal{G}_{\rho,\sigma}^{+}|G_{\rho,\sigma}^{+}}}
\arrow[d,phantom,"\simeq"{rotate=90}] \&\&
\inserted{\gamma_{\rho,\sigma}} \arrow[d,phantom,"\simeq"{rotate=90}]\&\\
1 \arrow[r] \& \mathcal{G}^{+\wedge} \arrow[r] \arrow[d,"\sim"] \&
M \arrow[r] \arrow[d,"\sim"] \& \widetilde{\Gamma}_{\mathbb{Q}}
\arrow[r] \arrow[d,phantom,"\simeq"{rotate=90}] \&1\\
(58)\quad1 \arrow[r] \& \mathcal{G}_{\rho,\sigma}^{+}
\arrow[r,"\unsure{\sim}"] \arrow[d,hook,"1"] \&
\mathcal{M}_{\rho,\sigma}^{\unsure{\natural|\sharp}}/\mu_{\rho} \arrow[r] \arrow[d,hook] \&
\Gamma_{\rho,\sigma}^{\unsure{\natural|\sharp}}/\mu_{\rho} \arrow[r] \arrow[d,hook] \&1\\
1 \arrow[r] \& \mathcal{G}_{\rho,\sigma}^{+} \arrow[r] \&
\mathcal{M}_{\rho,\sigma}/\mu_{\rho} \arrow[r] \& \Gamma_{\rho,\sigma}/\mu_{\rho} \arrow[r] \&1
\end{tikzcd}}

Ainsi, on trouve en passant que l'\unsure{hom.} canonique
\diagram{
\begin{tikzcd}[ampersand replacement=\&, column sep=small,row sep=large]
(59)\quad1 \arrow[r] \&
\underbrace{SZ_{\rho}\times SZ_{\sigma}}_{L_{\rho}\times L_{\sigma}} \arrow[r] \&
\Gamma_{\rho,\sigma} \arrow[r,"\inserted{\mathrm{\acute{e}pi}}"] \&
\gamma_{\rho,\sigma} \arrow[r] \&1\\
\& \text{admet une ``}\unsure{\text{bisection}}\text{''} \&\&
\deleted{\unsure{\Gamma_{\rho,\sigma}^{\natural}/\mu_{\rho}}} \arrow[u,"\sim"] \&
\Gamma_{\rho,\sigma}^{\unsure{\natural|\sharp}}\text{ (dont}
\end{tikzcd}}
l'intersection avec $\unsure{SZ_{\rho}}\times SZ_{\sigma}$ est $\mu_{\rho}\times1$)

\mspage{648}

\scanpage{650}

\[
(60)\qquad
\left\{\begin{aligned}
\Gamma_{\rho,\sigma}&\simeq
\Gamma_{\rho,\sigma}^{\unsure{\natural|\sharp}}\cdot
\underbrace{S\Gamma_{\rho,\sigma}}_{SZ_{\rho}\times SZ_{\sigma}=L_{\rho}\times L_{\sigma}}
\ ;\quad
\Gamma_{\rho,\sigma}^{\unsure{\natural|\sharp}}\cap S\Gamma_{\rho,\sigma}=\mu_{\rho}
\\
\Gamma_{\rho,\sigma}/\mu_{\rho}&\simeq
\left(\underbrace{\Gamma_{\rho,\sigma}^{\unsure{\natural|\sharp}}/\mu_{\rho}}
_{\gamma_{\rho,\sigma}\simeq\unsure{\widetilde{\Gamma}_{\mathbb{Q}}|\widetilde{\Pi}_{\mathbb{Q}}}}\right)
\cdot\left(L_{\rho}/\mu_{\rho}\times L_{\sigma}\right).
\end{aligned}\right.
\]

Pour dire ce que deviennent les sous-groupes remarquables
\[
M(0),\quad
\begin{array}{c}
M[0],\\
\inserted{\Vert}\\[-4pt]
\inserted{M^{\mathcal{G}}[0]}\\
\inserted{\Vert}\\[-4pt]
\inserted{M(\unsure{\omega|0}).L_{0}^{\mathcal{G}}}
\end{array}
\quad
\begin{array}{c}
M_{\unsure{a}}^{!}[0],\\
\inserted{\Vert}\\[-4pt]
\inserted{M^{\unsure{\Pi}}[0]}\\
\inserted{\Vert}\\[-4pt]
\inserted{M(\unsure{\omega|0}).L_{0}^{\Pi}}
\end{array}
\quad
\begin{array}{c}
M_{0}^{\unsure{!}}[0]\\
\inserted{\Vert}\\[-4pt]
\inserted{M^{\Pi_{0}}\unsure{[0]}}\\
\inserted{\Vert}\\[-4pt]
\deleted{\unsure{M(\omega)L_{0}|M(0)L_{0}}}\\
\inserted{M(\unsure{\omega|0})L_{0}^{\Pi_{0}}}
\end{array}
\]
\[
\begin{array}{c}
\unsure{\{M(\tfrac12)\}|M(\tfrac12)}\\
\inserted{\Vert}\\[-4pt]
\inserted{Z_{\sigma0}^{!}}
\end{array},\quad
M'(-\bar{\jmath}),\quad M'''(-\bar{\jmath}),\quad N_{\rho}^{?}
\]
de $M$ par l'homomorphisme
\diagram{
\begin{tikzcd}[ampersand replacement=\&, column sep=large]
M \arrow[r,"\sim"] \&
\mathcal{M}_{\rho,\sigma}^{\unsure{\natural|\sharp}}/\mu_{\rho} \arrow[r,hook] \&
\mathcal{M}_{\rho,\sigma},
\end{tikzcd}}
il faudrait d'abord introduire les groupes adéquats dans les $\mathcal{M}_{\rho,\sigma}$
généraux. Nous allons donc revenir sur ces groupes maintenant.

\mspage{649}

\scanpage{651}
\mspage{628}

\uline{VI}\quad Les groupes $\mathcal{M}_{\rho,\sigma}$ généraux:
récapitulation, \unsure{raffinements|raffinement}.

$1^\circ)$ Données. Ce sont celles de \uline{\unsure{V|IV}} $1^\circ)$ (p.~577).
En plus de $\mathcal{G}$ \inserted{(profini)} et de
\diagram{\begin{tikzcd}[ampersand replacement=\&]
\mathrm{Gl}(2,\mathbb{Z})^{\wedge} \arrow[r] \& \mathcal{G}
\end{tikzcd}}
hom continu surjectif, i.e. des générateurs $\rho,\sigma,\tau$
(ou $\rho,\sigma,\tau'$) de $\mathcal{G}$, on se donne aussi un sous-groupe ouvert
\[
(1)\qquad \mathcal{G}_{\unsure{0|a}}^{\unsure{\natural|\sharp}}
\subset\mathcal{G}^{+},\quad
\text{invariant dans }\mathcal{G}^{\unsure{*}}.
\]

Dans le cas universel $\mathcal{G}=\mathrm{Gl}(2,\mathbb{Z})^{\wedge}$,
la donnée de $\mathcal{G}_{\unsure{0|a}}^{\unsure{\natural|\sharp}}$
peut correspondre au point de vue où on regarde les sous-groupes
$M^{\unsure{\natural|\sharp}}$ de $M$
\deleted{qui}\inserted{\unsure{formé|formés} des $u$ qui}
\unsure{normalisent}\deleted{\unsure{ateur de|ent le}}
$\mathcal{G}^{\unsure{\natural|\sharp}}$
\inserted{et qui opèrent \unsure{trivialt} dans
$\mathcal{G}^{+}/\mathcal{G}^{\unsure{\natural|\sharp}}$ --- c'est}
\deleted{--- qui est \unsure{néct.} ouvert, et définit}
\inserted{un sous-groupe ouvert, qui définit}
dans $\unsure{\widetilde{\Gamma}_{\mathbb{Q}}|\widetilde{\Pi}_{\mathbb{Q}}}=M/\mathcal{G}^{+}$
un sous-groupe image
$\unsure{\widetilde{\Gamma}_{\mathbb{Q}}|\widetilde{\Pi}_{\mathbb{Q}}}^{\unsure{\natural|\sharp}}$
ouvert, dont on peut se proposer d'étudier les représentations dans des groupes
profinis convenables. Un cas intéressant sera celui où
\[
\mathcal{G}^{\unsure{\natural|\sharp}}=\Pi_{0},\quad
\text{donc }M^{\unsure{\natural|\sharp}}=M^{\prime\prime!},
\]
ou
\[
\mathcal{G}^{\unsure{\natural|\sharp}}=\Pi,\quad
\text{donc }M^{\unsure{\natural|\sharp}}=M^{!}.
\]
Si
\[
\mathcal{G}^{\unsure{\natural|\sharp}}=\mathcal{G},\quad
\text{donc }M^{\unsure{\natural|\sharp}}=M,
\]
on retrouve les constructions précédentes.
\deleted{\unsure{On peut}}

\mspage{650}

\scanpage{652}

Un cas intéressant aussi est celui où $\unsure{\mathcal{G}}$ est le groupe
des pts adéliques \inserted{entiers} d'un schéma en groupes sur $\mathbb{Z}$
(dans ce cas, il est vrai que $\rho$ et $\sigma,\tau$ n'engendrent
pas alors $\mathcal{G}^{\unsure{+|*}}$ que modulo groupe fini
i.e. ils engendrent \unsure{peut être} seulement un s-groupe ouvert.
Dans ce cas, il faudrait plutôt \uline{\unsure{paraphraser}} les constructions
qui suivent dans le cas schématique
J'y reviendrai par la suite. Un choix intéressant d'un
$\mathcal{G}^{\unsure{\natural|\sharp}}$ est alors la
\uline{composante neutre} de $\mathcal{G}$.

J'hésite à introduire ces notations
$Z_{\rho,\sigma}^{\unsure{\natural|\sharp}},\mathcal{M}_{\rho,\sigma}^{\unsure{\natural|\sharp}}$
etc --- pour ne pas surcharger des notations déjà lourdes. Donc je préfère
utiliser les notations plus simples $Z_{\rho,\sigma}$, etc --- étant entendu
que dans ces constructions, le groupe $\mathcal{G}^{\unsure{\natural|\sharp}}$
intervient --- quitte à réintroduire la notation $\unsure{\natural|\sharp}$
en cas de besoin.
On pose
\[
(3)\qquad
\mathcal{G}_{\tau}^{\unsure{\natural|\sharp}}
=\overbrace{\{1,\tau\}}^{\mu_{\tau}}\,.
\mathcal{G}^{\unsure{\natural|\sharp}}\subset\mathcal{G}.
\]

\mspage{651}

\scanpage{653}
\mspage{629}

$2^\circ)$ \quad $Z_{\rho,\sigma},\ \gamma_{\rho,\sigma}$.

Soit
\[
(4)\qquad
\begin{array}{l}
 Z_{\rho,\sigma}\subset
 \operatorname{Aut}(\mathcal{G}^{+})\times
 \widehat{\mathbb{Z}}^{*}\times\mu\times\mu\\
 \phantom{Z_{\rho,\sigma}}\Vert\,\mathrm{def}\\
 \left\{(u,\mu,\varepsilon_2,\varepsilon_3)\ \left|
 \begin{array}{ll}
 \multicolumn{2}{l}{\deleted{\unsure{u\ \text{normalise}\ \mathcal{G}^{\natural}\ 
 \text{et induit l'identité dans}\ \mathcal{G}^{+}/\mathcal{G}^{\natural}}}}\\
 a)&\exists\alpha\in\mathcal{G}_{\tau}^{\unsure{\natural|\sharp}},
 \ \text{avec }\varepsilon_{\mathcal{G}}(\alpha)=\varepsilon_3,
 \quad u(\rho)=\alpha(\rho)\\
 b)&\exists\beta\in\mathcal{G}_{\tau}^{\unsure{\natural|\sharp}},
 \ \text{avec }\varepsilon_{\mathcal{G}}(\beta)=\varepsilon_2,
 \quad u(\sigma)=\beta(\sigma)\\
 c)&u(\varepsilon_0)=\varepsilon_0^\mu\\
 d)&u\ \text{normalise }\mathcal{G}^{\unsure{\natural|\sharp}}
 \inserted{\ \text{\uline{et} induit l'identité sur }
 (\mathcal{G}^{+}/\mathcal{G}^{\unsure{\natural|\sharp}})/\operatorname{Im}\Pi}
 \end{array}\right.\right\}.
\end{array}
\]
\margin{\unsure{cf. page} \unsure{635} \unsure{suite}
\deleted{\unsure{une condition plus faible}}
\unsure{cette condition}.}
\margin{\deleted{La dernière condition \unsure{rajoutée} dans d)
(\unsure{en principe}), est \unsure{conséquence} de a)b) si
$\varepsilon_2\varepsilon_3=1$
(\unsure{auquel cas} $u|(\mathcal{G}^{+}/\mathcal{G}^{\unsure{\natural|\sharp}})$
\deleted{l'identité} \unsure{induite par} $\varepsilon[\tau]$,
\unsure{où} $\varepsilon=\varepsilon_2=\varepsilon_3\in\{\pm1\}$)}.}

\uline{Proposition.} \quad $Z_{\rho,\sigma}$ est un \uline{sous-groupe}
du groupe produit.

\uline{Dém.} Les conditions c), d) définissent \unsure{évidemment}
\deleted{un} \inserted{chacune} sous-groupe.
\unsure{Compte tenu que}
\[
\tau(\sigma)=-\unsure{\sigma},
\]
la condition b) signifie encore que
\[
\unsure{(5^{*})}\qquad
b')\quad u(\sigma)\ \text{est }\mathcal{G}^{\unsure{\natural|\sharp}}
\text{-conjugué à : }\
\deleted{\unsure{\varepsilon_2}}\sigma^{\varepsilon_2}
\bigl(\deleted{\unsure{\varepsilon_2\sigma}}\inserted{=\varepsilon_2\sigma}\bigr)
\]
et sous cette forme il est clair que la condition \unsure{c')|b')}
définit un sous-groupe. Je ne pense pas que la condition a)
définisse à elle seule un sous-groupe. Mais si
$\unsure{\mathfrak{g}_{(b)}|G_{(b)}}$
\deleted{$\unsure{\mathfrak{g}_{(a,b)}}$}
désigne le sous-\unsure{grpe} du produit défini par ces conditions
\deleted{\unsure{a)}}\inserted{b)}, on a \unsure{à montrer} que
$\unsure{\mathfrak{g}|G}_{\deleted{\unsure{(a,b,c)}}\inserted{(a,b)}}$
est un sous-groupe du \deleted{produit}
$\unsure{\mathfrak{g}|G}_{\deleted{\unsure{(a,b)}}\inserted{(b)}}$
--- ce qui prouvera la prop.
Soient $\uline{u},\uline{u}'\in\unsure{\mathfrak{g}_{(b)}|G_{(b)}}$,
correspondant à
\inserted{$(u,\mu,\varepsilon_2,\varepsilon_3)$ et
$(u',\mu',\varepsilon'_2,\varepsilon'_3)$, et à}
$\alpha,\alpha'\in\mathcal{G}^{\unsure{\natural|\sharp}}_\tau$,
\unsure{de sorte} que l'on a

\mspage{652}

\scanpage{654}

\[
(6)\qquad
u(\rho)=\alpha(\rho),\quad u'(\rho)=\alpha'(\rho),
\inserted{\quad\text{avec }\begin{array}{l}
\varepsilon_{\mathcal{G}}(\alpha)=\varepsilon_3,\\
\unsure{\alpha_{\mathcal{G}}|\varepsilon_{\mathcal{G}}}(\alpha')=\varepsilon'_3.
\end{array}}
\]
\unsure{Prouvons} que
$u'u\in\unsure{\mathfrak{g}_{(a)(b)}|G_{(a,b)}}$
--- i.e. qu'il satisfait aussi à \unsure{(b)|(a)}.
Si $\varepsilon_3=1$ i.e.
$\alpha\in\mathcal{G}^{\unsure{\natural|\sharp}}
\inserted{\subset\mathcal{G}^{+}}$,
$u'(\alpha)$ est défini et on a
\[
\begin{aligned}
(u'u)(\rho)&=u'(u(\rho))=u'(\alpha(\rho))
=u'(\alpha)(u'(\rho))=u'(\alpha)(\alpha'(\rho))\\
&=(u'(\alpha)\alpha')(\rho)=\alpha''(\rho),
\end{aligned}
\]
avec
\[
(7)\qquad
\begin{array}{l|l}
\alpha''=u'(\alpha)\alpha'\in
\mathcal{G}^{\unsure{\natural|\sharp}}_\tau&
\text{Cas }\varepsilon_3=1\\
\varepsilon_{\mathcal{G}}(\alpha'')
=\varepsilon_{\mathcal{G}}(\alpha')
=\varepsilon'_3=\varepsilon_3\varepsilon'_3&
\end{array}
\]
(puisque $u'(\alpha)\in\mathcal{G}^{\unsure{\natural|\sharp}}$
i.e. $\varepsilon_{\mathcal{G}}(u'(\alpha))=1$, et $\varepsilon_3=1$.

Dans ce cas on \unsure{gagne}. Si $\varepsilon_3=-1$, on écrit
\[
\alpha=\alpha_0\tau\qquad
\alpha_0\in\mathcal{G}^{\unsure{\natural|\sharp}},
\]
et on a
\[
u(\rho)=(\alpha_0\tau)(\rho)=\alpha_0(\tau(\rho))
=\alpha_0\bigl(\sigma^{\unsure{-1}}(\rho^{-1})\bigr)
=(\alpha_0\sigma)(\rho^{-1}),
\]
où maintenant $\alpha_0\sigma\in\mathcal{G}^{+}$, donc
$u'(\alpha_0\sigma)$ est défini, et on aura
\[
\unsure{(u'u)(\rho)}
=u'(u(\rho))
=u'(\alpha_0\sigma)
\bigl(\underbrace{u'(\rho^{-1})}_{\alpha'(\rho^{-1})}\bigr)
=\bigl(u'(\alpha_0\sigma)\alpha'\bigr)(\rho^{-1})
\]
et comme
\[
\deleted{\unsure{\rho^{-1}}}\qquad
\rho^{-1}=\sigma\tau(\rho)
\]
on trouve
\[
(8)\qquad
(u'u)(\rho)=
\deleted{\unsure{u'(\alpha_0\sigma)\alpha'\sigma\tau}}
\alpha''(\rho)
\]
avec
\[
\left\{
\begin{array}{l}
\alpha''=\varepsilon
\overbrace{u'(\alpha_0)}^{\inserted{\in\mathcal{G}^{\unsure{\natural|\sharp}}}}
\overbrace{u'(\sigma)}^{\inserted{\in\mathcal{G}^{+}}}
\overbrace{\alpha'}^{\inserted{\in\mathcal{G}^{\unsure{\natural|\sharp}}_\tau}}
\underset{\deleted{\unsure{\tau'}}}{\sigma\tau}
\qquad\text{où}\\
\varepsilon\in\mu=\{\pm1\}
\unsure{\hookrightarrow|\Rightarrow}\operatorname{Centre}\mathcal{G}^{+}
\end{array}
\right.
\]
\deleted{\unsure{est}} est choisi de façon que $\alpha''$ (si possible)
soit dans $\mathcal{G}^{\unsure{\natural|\sharp}}_\tau$.
Notons que tous les facteurs de ce produit sont dans
$\unsure{\mathcal{G}^{\natural}_\tau\mathcal{G}^{+}|
\mathcal{G}^{\natural}_\tau/\mathcal{G}^{+}}$,
\unsure{\uline{en}} tous cas $\alpha''\in\mathcal{G}$.

\mspage{653}

\scanpage{655}
\mspage{630}

Réduisons mod $\mathcal{G}^{\unsure{\natural|\sharp}}$,
la relation $\alpha''\in\mathcal{G}^{\unsure{\natural|\sharp}}_\tau$
s'écrit
\[
\dot{\alpha}''\in\{1,\dot{\tau}\}\quad\text{i.e.}\quad
\dot{\varepsilon}\,u'(\sigma)^{\cdot}\dot{\alpha}'\dot{\sigma}\dot{\tau}
\in\{1,\dot{\tau}\}
\quad\text{soit}\quad
\dot{\varepsilon}\,u'(\sigma)^{\cdot}\dot{\alpha}'\dot{\sigma}
\in\{1,\dot{\tau}\}.
\]
\deleted{\unsure{on aurait}, puisque $\sigma\tau=-\tau\sigma$}
\inserted{On distingue \unsure{encore} deux cas}

$1^\circ)$ \quad $\varepsilon'_3=1$ i.e. $\dot{\alpha}'=1$
\deleted{$\unsure{\dot{\varepsilon}\,u'(\sigma)^{\cdot}
\dot{\alpha}'\dot{\tau}}$}
la relation s'écrit
\[
(*)\qquad \dot{\varepsilon}\,u'(\sigma)^{\cdot}\dot{\sigma}=1
\]
$2^\circ)$ \quad $\varepsilon'_3=-1$ i.e. $\dot{\alpha}'=\dot{\tau}$,
comme $\tau\sigma=-\sigma\tau$, la relation s'écrit
\[
(**)\qquad -\dot{\varepsilon}\,u'(\sigma)^{\cdot}\dot{\sigma}=1,
\]
donc dans tous les cas, la relation s'écrit
\[
(9)\qquad
\left\{
\begin{array}{l}
u'(\sigma)^{\cdot}=(\varepsilon\varepsilon'_3)^{\cdot}
\dot{\sigma}^{-1}
\quad\deleted{\unsure{\text{ou}}}\ \text{soit encore}\\
u'(\sigma)^{\cdot}=(-\varepsilon\varepsilon'_3)^{\cdot}\dot{\sigma}
\end{array}
\right.
\]
Or \deleted{utilisons l'hyp. b) sur $u'$,}
\inserted{supposons maintenant qu'on a b) pour $u'$,}
qui \deleted{\unsure{signifie}}\inserted{implique}
\[
u'(\sigma)^{\cdot}
=\deleted{\unsure{\varepsilon_2'}}
\inserted{\unsure{(\varepsilon'_2)^{\cdot}}}\dot{\sigma},
\]
on voit donc qu'il suffit de choisir $\varepsilon$ de façon que
$-\varepsilon\varepsilon'_3=\unsure{\dot{\varepsilon}'_2|\varepsilon'_2}$
i.e.
\[
\varepsilon=-\varepsilon'_2\varepsilon'_3
\]
donc on aura
\[
\deleted{\unsure{(10)}}\qquad
\alpha''=-\varepsilon'_2\varepsilon'_3\,
u'\bigl(
\underbrace{\alpha_0\deleted{\unsure{\tau}}\sigma}
_{\inserted{\alpha\tau\sigma=\alpha\tau'}}
\bigr)\alpha'\underbrace{\sigma\tau}_{\inserted{-\tau'}}
\]
i.e.
\[
(10)\qquad
\alpha''=\varepsilon'_2\varepsilon'_3\,u'(\alpha\tau')
(\alpha'\tau')
\quad\big|\quad\text{Cas }\varepsilon_3=-1
\]
formule qui a un sens, car $\alpha\tau'\in\mathcal{G}^{+}$ donc
$u'(\alpha\tau')$ est défini, et l'argument précédent prouve que
l'élément en question \deleted{\unsure{est dans}}\inserted{satisfait}
\[
(11)\qquad
\alpha''\in\mathcal{G}^{\unsure{\natural|\sharp}}_\tau,\quad
\varepsilon_{\mathcal{G}}(\alpha'')
=\underbrace{\varepsilon_{\mathcal{G}}(\alpha)}_{-1}
\varepsilon_{\mathcal{G}}(\alpha')
\]
en distinguant les cas \unsure{(*)} et $(**)$ de (9).
$\overline{\mathrm m}$ conclusion dans le cas premier où
$\varepsilon_{\mathcal{G}}(\alpha)=1$ (7).
\inserted{Cela achève la démonstration.}

\mspage{654}

\scanpage{656}

\uline{Remarques} \deleted{a) Pour que la condition b) soit
\unsure{multiplicative}, on n'a pas besoin de la condition a) dans
\unsure{tte} sa force, il suffit que $u$ dans
$\mathcal{G}^{+}/\mathcal{G}^{\unsure{\natural|\sharp}}$
normalise l'image de $L_\sigma$, i.e. qu'on ait
\[
u(\sigma)\equiv\pm\sigma
\quad(\mathcal{G}^{\unsure{\natural|\sharp}}).
\]}

\deleted{b)} Si $\mathcal{G}$ est “assez gros”, l'hom de projection
est injectif

(12)\quad
\diagram{\begin{tikzcd}[ampersand replacement=\&, column sep=large]
Z_{\rho,\sigma}\arrow[r,hook]\&\operatorname{Aut}(\mathcal{G}^{+})
\end{tikzcd}}
\quad (cf. p. 578)

\uline{Exemples} dans le cas universel

a) \quad $\mathcal{G}^{\unsure{\natural|\sharp}}=\mathcal{G}^{+}
=\deleted{\unsure{\mathcal{G}}}\operatorname{SL}(2,\mathbb{Z})^{\wedge}$.
\deleted{On a ici la} \unsure{Condition|condition} (4) d) est vide,
les conditions a), b) équivalent respectivement au fait que
$u(\rho)$ conjugué à $\rho^{\varepsilon_3}$,
$u(\sigma)$ conjugué à $\sigma^{\varepsilon_2}$,
i.e. \unsure{retombe} sur la définition (6) p. 578, qui donne
\[
\deleted{Z_{\rho,\sigma}=M(\unsure{\omega|0}).L_0^{\mathcal{G}}
=M^{\mathcal{G}}[0].}
\]
\[
(13)\qquad
Z_{\rho,\sigma}=M(0).L_0^{\mathcal{G}}
=M^{\mathcal{G}}[0]=M[0]
\]

\margin{\uline{NB} A priori on a plutôt
$Z_{\rho,\sigma}\subset\unsure{\widetilde{M}_{0,3}|M_{0,3}}$, mais
\diagram{\begin{tikzcd}[ampersand replacement=\&, column sep=large,row sep=large]
M[0]\arrow[r,"\sim"]\arrow[d,equal]\&
\widetilde{M}_{0,3}[0]\arrow[d,equal]\\
M(\unsure{\omega|0}).L_0^{\mathcal{G}}\arrow[r,"\sim"]\&
\unsure{M_{0,3}(0)}.
L_0^{\tikz[baseline=(g.base)]{\node[draw,ellipse,inner sep=1pt]
(g){$\scriptstyle\unsure{\mathcal{G}_{0,3}}$};}}
\end{tikzcd}}}

b) \quad
$\mathcal{G}^{\unsure{\natural|\sharp}}=\Pi=\Pi_0\times\mu$.
On trouve maintenant
\[
(14)\qquad
Z_{\rho,\sigma}=M(0).L_0^\Pi
=M^\Pi[0]=\unsure{M^{!}[0]|M^{\prime!}[0]}
\]
(noté $Z_{\rho,\sigma}^{!}$ dans $\mathrm{V}$ n$^\circ$ 8).

\margin{Dans ce cas, il est immédiat que $Z_{\rho,\sigma}$
opère \unsure{trivial\textsuperscript{mt}} même sur
$\mathcal{G}^{+}/\mathcal{G}^{\unsure{\natural|\sharp}}_\tau$,
car il en est ainsi du
\deleted{\unsure{NB} \unsure{tous les sous-groupes}
\unsure{contenus par} a), b) de (4),
$\unsure{\mathcal{G}^{\natural}=M(0)_2.L_0^\Pi|
\mathcal{G}^{\natural}=M(\omega)_2.L_0^\Pi}$}.}
\margin{$\unsure{M^{!}[0]|M^{\prime!}[0]}=M^\Pi[0]$.}
\margin{\unsure{expliciter} à nouveau pourquoi
\unsure{(i.e. pourquoi)} que
$(\varepsilon_0,1,1,1)\unsure{\notin|\in}Z_{\rho,\sigma}$}

c) \quad $\mathcal{G}^{\unsure{\natural|\sharp}}=\Pi_0$.
Les conditions a) et b) \inserted{de (4)} n'ont pas changé
\unsure{p.r.v.|p.r. au}
\deleted{\unsure{avant}}\inserted{\uline{cas précédent}},
puisque $\Pi=\Pi_0\times\mu$ et $\mu$ est central, par contre la

\mspage{655}

\scanpage{657}
\mspage{631}

condition d) devient nettement plus stricte, on trouve p.ex. que
$Z_{\rho,\sigma}\cap\mathcal{G}^{+}=L_0^{\Pi_0}$
(d'indice 2 dans le groupe correspondant $L_0^\Pi$ du cas précédent),
et comme l'opération de $M(0)$ \inserted{$\subset M^!$} sur
$\mathcal{G}^{+}/\Pi_0$ se fait par le \unsure{caractère}
$\varepsilon_2^{\unsure{\#|\prime}}$, on voit que l'on a maintenant
\[
(15)\qquad
Z_{\rho,\sigma}=M''(0).L_0^{\Pi_0}=M''[0]_0
\]
\[
\left\{
\begin{array}{l}
M''(0)=\operatorname{Ker}(M(0)\xrightarrow{\varepsilon_2}\mu)\\
M''[0]_0=\operatorname{Ker}(M[0]_0\xrightarrow{\varepsilon_2}\mu)\\
M[0]_0=M(\unsure{\omega|0}).L_0^{\Pi_0}
\end{array}
\right.
\]

\margin{\deleted{Il en est ainsi si dans la définition (4) de
$Z_{\rho,\sigma}$ on exige dans d) que $u$
\unsure{induise|produise} l'op. triviale sur
$\mathcal{G}^{+}/\mathcal{G}^{\unsure{\natural|\sharp}}$,
condition que j'ai fini par abandonner comme inutile.
\unsure{Dans la nouvelle version} les cas
$\mathcal{G}^{\unsure{\natural|\sharp}}=\Pi$ et
$\mathcal{G}^{\unsure{\natural|\sharp}}=\Pi_0$
donnent le $\overline{\mathrm m}$ $Z_{\rho,\sigma}$ (15).
Mais je voudrais obtenir
\[
M^{\Pi_0}[0]=\unsure{M[0]^!_0|M[0]_0}
=M(\unsure{\omega|0}).\Pi_0
\]
--- le groupe de b) est \unsure{trop} grand, celui de c)
\unsure{un peu tr.}
\uline{\unsure{(ancienne version)}}
\inserted{un peu trop petit!}}}

\margin{\deleted{Mais avec la nouvelle définition de
$Z_{\rho,\sigma}$ (condition ajoutée à (4)) on aura finalement
\[
Z_{\rho,\sigma}=M^![0]
\]
comme dans le cas b) précédent.
\inserted{On aurait pu \unsure{introduire}
un $\unsure{M^![0]_0|M^{\prime!}[0]_0}
=M(\unsure{0|\omega}).\unsure{L_0^{\Pi_0}|\Pi_0}$.}}}

\margin{Cette remarque sur un
$\mathcal{G}^{\unsure{\natural|\sharp}}_0$
devient maintenant inutile.}

Ces deux \inserted{derniers} exemples suggèrent l'opportunité,
dans le cas \inserted{général} où on disposerait d'un ss-groupe
\deleted{\unsure{$Z$}} $\mathcal{G}^{\unsure{\natural|\sharp}}_0$
de $\mathcal{G}^{\unsure{\natural|\sharp}}$,
\inserted{invariant dans $\mathcal{G}$}, tel qu'on ait
\[
(16)\qquad
\mathcal{G}^{\unsure{\natural|\sharp}}
=\mathcal{G}^{\unsure{\natural|\sharp}}_0\times\mu
\quad\deleted{\unsure{\subset|\Leftarrow}}
\]
\inserted{(ici $\mu$ est le sous-groupe central universel
$\mu=\{1,\rho^3\}=\{1,\sigma^2\}$)}
de définir $Z_{\rho,\sigma}$ en termes de
$\mathcal{G}^{\unsure{\natural|\sharp}}$, et de
\uline{choisir} les $\alpha,\beta$ de la définition (4) dans
\deleted{$\unsure{\mathcal{G}^{\natural}_{0\tau}}$}
\[
(17)\qquad
\mathcal{G}^{\unsure{\natural|\sharp}}_{0\tau}
\overset{\mathrm{def}}{=}
\mu_\tau.\mathcal{G}^{\unsure{\natural|\sharp}}_0.
\]

\uline{Remarque} On est tenté de remplacer la déf (4)
\inserted{(pour a) et b))}
par la définition apparemment plus simple de la p. 578,
qui rend évident qu'on a un sous-groupe.
Mais elles ne \deleted{\unsure{sont compatibles}}
\inserted{donnent pas, $\overline{\mathrm m}$ si
$\mathcal{G}^{\unsure{\natural|\sharp}}\subset\Pi$
(\uline{cas universel})}
\deleted{avec \unsure{si on suppose}}
une action triviale de
\deleted{$\mathcal{G}^{+}/\mathcal{G}^{\unsure{\natural|\sharp}}
\longrightarrow\mathcal{G}_3$}
\inserted{$Z$ sur $\mathcal{G}^{+}/\Pi$,
comme on le \deleted{\unsure{voit}}
\inserted{voit} immédiatement,
\deleted{\unsure{$\rho,\sigma$}}}

\mspage{656}

\scanpage{658}

qu'\unsure{en} se \unsure{restreignant} à des, $(u,\mu,\varepsilon_2,\varepsilon_3)$ avec $\varepsilon_3=1$,
car $\dot{\rho}^{-1}$ n'est pas égal à $\dot{\rho}$ ! Il est vrai \deleted{que}
(dans le cas universel p. ex.) que $u(\rho)$ est conjugué \uline{dans $\mathcal{G}^{+}$}
à $\rho^{\varepsilon_3}$, mais en général il ne le sera pas par un élément (disons) de $\Pi$.

On pose
\[
(18)\qquad
L_0^{\unsure{\natural|\sharp}}
=\deleted{\unsure{L_0}\cap\mathcal{G}^{\unsure{\natural|\sharp}}}
\inserted{\bigl\{l\in L_0^{\mathcal{G}}\mid
(\operatorname{int}(l),1,1,1)\in Z_{\rho,\sigma}\bigr\}}
\qquad
\left(\text{où }L_0^{\mathcal{G}}\subset\mathcal{G}^{+}
\text{ est l'\unsure{ens} des }\varepsilon_0^n,\ n\in\widehat{\mathbb{Z}}\right)
\]
\margin{\uline{NB} On a
$L_0^{\unsure{\natural|\sharp}}\supset L_0\cap\mathcal{G}^{\unsure{\natural|\sharp}}$,
et l'inclusion peut être stricte, cf. exemple c) ci-dessus\unsure{\ldots}
où \uline{$L_0^{\unsure{\natural|\sharp}}=L_0^\Pi$},
$L_0\cap\mathcal{G}^{\unsure{\natural|\sharp}}=L_0^{\Pi_0}$.
\unsure{ce qui} \unsure{surprend}.}

et on a un hom canonique
\diagram{
\begin{tikzcd}[ampersand replacement=\&]
(19)\quad L_0^{\unsure{\natural|\sharp}}
  \arrow[r,"\unsure{i}"] \& Z_{\rho,\sigma},\\
l \arrow[r,mapsto] \& (\operatorname{int}_{\mathcal{G}^{+}}(l),1,1,1)
\end{tikzcd}
}
injectif ssi on a
\[
(20)\qquad
L_0^{\unsure{\natural|\sharp}}\cap\operatorname{Centr}(\mathcal{G}^{+})=\{1\}
\]
l'image invariante. On pose
\[
(21)\qquad
\gamma_{\rho,\sigma}=Z_{\rho,\sigma}/L_0^{\unsure{\natural|\sharp}},
\]
d'où
\diagram{
\begin{tikzcd}[ampersand replacement=\&]
(22)\quad 1 \arrow[r] \&
L_0^{\unsure{\natural|\sharp}} \arrow[r] \&
Z_{\rho,\sigma} \arrow[r,"\delta_Z"] \&
\gamma_{\rho,\sigma} \arrow[r] \& 1,
\end{tikzcd}
}
\diagram{
\begin{tikzcd}[ampersand replacement=\&]
 \& \& \mu\\
(23)\quad Z_{\rho,\sigma} \arrow[r] \&
\gamma_{\rho,\sigma}
  \arrow[ur,"\varepsilon_{3\gamma}"]
  \arrow[r,"\chi_\gamma"]
  \arrow[dr,"\varepsilon_{2\gamma}"'] \&
\widehat{\mathbb{Z}}^{*}\\
 \& \& \mu
\end{tikzcd}
}
\unsure{)} \unsure{déduites} des
\diagram{
\begin{tikzcd}[ampersand replacement=\&, row sep=0pt]
 \& \varepsilon_3\\
(u,\mu,\varepsilon_2,\varepsilon_3)
  \arrow[ur,mapsto] \arrow[r,mapsto] \arrow[dr,mapsto]
  \& \mu\\
 \& \varepsilon_2.
\end{tikzcd}
}

\deleted{
\uline{Remarques} On \deleted{désigne} pose,
\[
(24)\qquad
\left\{
\begin{array}{lll}
\gamma'_{\rho,\sigma}=\operatorname{Ker}\varepsilon_{3\gamma},&
\gamma''_{\rho,\sigma}=\operatorname{Ker}\unsure{\varepsilon_{2,\sigma}|\varepsilon_{2\gamma}},&
\gamma'''_{\rho,\sigma}=\operatorname{Ker}\varepsilon_{2\gamma}\varepsilon_{3\gamma}\\
Z'_{\rho,\sigma}=\ldots&
Z''_{\rho,\sigma}=\ldots&
Z'''_{\rho,\sigma}=\ldots
\end{array}
\right.
\qquad
(\text{images inverses dans }Z_{\rho,\sigma}\unsure{)}
\]
}
\mspage{657}

\scanpage{659}
\mspage{632}

On pose
\[
(24)\qquad
\left\{
\begin{array}{l}
\gamma'_{\rho,\sigma}=\operatorname{Ker}\varepsilon_{3\gamma},\
\gamma''_{\rho,\sigma}=\operatorname{Ker}\varepsilon_{2\gamma},\
\gamma'''_{\rho,\sigma}=\operatorname{Ker}\varepsilon_{2\gamma}\varepsilon_{3\gamma},\
\gamma^\circ_{\rho,\sigma}=\gamma'\cap\gamma''=\gamma''\cap\gamma'''
=\gamma'''\cap\gamma'\\[4pt]
Z'_{\rho,\sigma},Z''_{\rho,\sigma},Z'''_{\rho,\sigma},Z^\circ_{\rho,\sigma}
\text{ sont les images inv. des }\deleted{\unsure{Z_\rho}}
\gamma'_{\rho,\sigma}\text{ etc}\\
\hfill\text{dans }Z_{\rho,\sigma}
\end{array}
\right.
\]
Ce sont donc des sous groupes d'indice $1$, \unsure{ou} $2$ ou quatre de
$\gamma_{\rho,\sigma}$ \unsure{respt} $Z_{\rho,\sigma}$ --- l'indice $4$ ne pouvant
se présenter que pour $\gamma^\circ_{\rho,\sigma},Z^\circ_{\rho,\sigma}$.

\uline{Remarques} 1) On a donc
\[
(25)\qquad
Z^\circ_{\rho,\sigma}\simeq
\left\{
(u,\mu)\in\operatorname{Aut}(\mathcal{G}^{+})\times\widehat{\mathbb{Z}}^{*}
\left|
\begin{array}{ll}
\mathrm{a})&u(\rho)\text{ est }\mathcal{G}^{\unsure{\natural|\sharp}}\text{-conj. à }\rho\\
\mathrm{b})&u(\sigma)\quad\uline{\hspace{4em}}\quad\sigma\\
\mathrm{c})&u(\varepsilon_0)=\varepsilon_0^\mu
\end{array}
\right.
\right\}\unsure{,}
\]
\margin{\unsure{Plus} généralement,
\[
Z'''_{\rho,\sigma}\simeq
\left\{
(u,\mu,\varepsilon)\in\operatorname{Aut}(\mathcal{G}^{+})
\times\widehat{\mathbb{Z}}^{*}\times\mu
\left|
\begin{array}{ll}
\mathrm{a})&
u(\rho)\ \deleted{\unsure{\text{est }\varepsilon[\tau](\rho)}}
\inserted{\mathcal{G}^{\unsure{\natural|\sharp}}\text{-conj.}}
\text{ à }\varepsilon[\tau](\rho)\\
\mathrm{b})&
u(\sigma)\text{ est }\mathcal{G}^{\unsure{\natural|\sharp}}\text{-conj. à }
\varepsilon[\tau](\sigma)\\
\mathrm{c})&u(\varepsilon_0)=\varepsilon_0^\mu
\end{array}
\right.
\right\}\unsure{,}
\]}
(isomorphisme de \deleted{sous} groupes). La condition d) de (4) ---
stabilité de $\mathcal{G}^{\unsure{\natural|\sharp}}$ sous $u$ --- est conséquence
ici de a) b) c), car on \unsure{aura}
\[
\begin{array}{ll}
u(\rho)\equiv\rho&(\mathcal{G}^{\unsure{\natural|\sharp}})\\
u(\sigma)\equiv\sigma&(\mathcal{G}^{\unsure{\natural|\sharp}})
\end{array}
\]
ce qui implique
\[
(25\mathrm{bis})\qquad
u(x)\equiv x\quad(\mathcal{G}^{\unsure{\natural|\sharp}})
\qquad \forall x\in\mathcal{G}^{+},
\]
i.e. la stabilité de $\mathcal{G}^{\unsure{\natural|\sharp}}$ \uline{et} une action triviale de
$u$ sur $\mathcal{G}/\mathcal{G}^{\unsure{\natural|\sharp}}$.

2) On a encore l'élément remarquable
\[
(26)\qquad
\tau_Z\in Z_{\rho,\sigma},\qquad \tau_Z=(\tau,-1,-1,-1)
\]
\margin{C'est une raison de plus pour ne pas exiger \unsure{ds.} (4) d) que $u$ opère
trivialement sur $\mathcal{G}^{+}/\mathcal{G}^{\unsure{\natural|\sharp}}$,
\unsure{ce qui} \unsure{me ferait|on ferait} perdre $\tau$ !}

$3^\circ)$\quad $S_0G_{\rho,\sigma},G_{\rho,\sigma},
S_0G_{\rho,\sigma}^{\unsure{\natural|\sharp}},G_{\rho,\sigma}^{\unsure{\natural|\sharp}}$.

Procédant comme page \unsure{577|579}, on trouve
\diagram{
\begin{tikzcd}[ampersand replacement=\&]
Z_{\rho,\sigma}\arrow[r]\&\operatorname{Aut}(\mathcal{G}^{+})
\end{tikzcd}
}
d'où des produits $1/2$ directs
\[
(27)\qquad
S_0G_{\rho,\sigma}\overset{\mathrm{dfn}}=
Z_{\rho,\sigma}\mathbin{.}\mathcal{G}^{+},\supset
S_0G_{\rho,\sigma}^{\unsure{\natural|\sharp}}\overset{\mathrm{dfn}}=
Z_{\rho,\sigma}\mathbin{.}\mathcal{G}^{\unsure{\natural|\sharp}}
\]
et un hom
\diagram{
\begin{tikzcd}[ampersand replacement=\&]
(28)\quad L_0^{\unsure{\natural|\sharp}}
  \arrow[r,hook,"\unsure{i_{L_0,S_0G}|i_{L_0,G}}"]\&
S_0G_{\rho,\sigma}^{\unsure{\natural|\sharp}}
  \arrow[r,phantom,"\subset"]\&S_0G_{\rho,\sigma}\\
l\arrow[r,mapsto]\&
\deleted{\unsure{\iota(l)|i(l)}}\,i_{L_0,Z}(l)\mathbin{.}l^{-1}
\end{tikzcd}
}
\mspage{658}

\scanpage{660}

d'image invariante, ce qui permet de poser
\[
(29)\qquad
G_{\rho,\sigma}\overset{\mathrm{dfn}}=
\underbrace{S_0G_{\rho,\sigma}}_{Z_{\rho,\sigma}\mathbin{.}\mathcal{G}^{+}}
/\operatorname{Im}L_0^{\unsure{\natural|\sharp}}
\ ,\supset\
G_{\rho,\sigma}^{\unsure{\natural|\sharp}}
=S_0G_{\rho,\sigma}^{\unsure{\natural|\sharp}}
/\operatorname{Im}\unsure{L_0^{\natural}|L_0^{\sharp}}
\]
d'où suites exactes
\diagram{
\begin{tikzcd}[ampersand replacement=\&]
(30)\quad 1\arrow[r,dash]\&
L_0^{\unsure{\natural|\sharp}}\arrow[r]\arrow[d,equal]\&
S_0G_{\rho,\sigma}^{\unsure{\natural|\sharp}}\arrow[r]\arrow[d,hook]\&
G_{\rho,\sigma}^{\unsure{\natural|\sharp}}\arrow[r,dash]\arrow[d,hook]\&1,\\
1\arrow[r,dash]\&
L_0^{\unsure{\natural|\sharp}}\arrow[r,dash]\&
S_0G_{\rho,\sigma}\arrow[r]\&G_{\rho,\sigma}\arrow[r,dash]\&1
\end{tikzcd}
}
et diagramme commutatif
\diagram{
\begin{tikzcd}[ampersand replacement=\&, column sep=small,row sep=small]
\&Z_{\rho,\sigma}
 \arrow[dr,hook,"i_Z"]
 \arrow[ddr,"{\deleted{\unsure{\longrightarrow}}}" description]
\&\&\\
(31)\quad L_0^{\unsure{\natural|\sharp}}
 \arrow[ur,hook,"i_{L_0,Z}"]
 \arrow[ddr,hook,"i_{L_0,\mathcal{G}}"']
\&\&G_{\rho,\sigma}^{\unsure{\natural|\sharp}}
 \arrow[dr,hook]\&\\
\&\&\deleted{\unsure{G_{\rho,\sigma}}}\&G_{\rho,\sigma}\\
\&\mathcal{G}^{\unsure{\natural|\sharp}}
 \arrow[uur,hook,"i_{\mathcal{G}^{\unsure{\natural|\sharp}}}"]
 \arrow[ur,"{\deleted{\unsure{\longrightarrow}}}" description]
 \arrow[dr,hook]\&\&\\
\&\&\mathcal{G}^{+}\arrow[uur,hook,"i_{\mathcal{G}}"']\&
\end{tikzcd}
}
s'insérant dans diagramme commutatif de suites exactes
\diagram{
\begin{tikzcd}[ampersand replacement=\&]
(32)\quad 1\arrow[r,dash]\&
L_0^{\unsure{\natural|\sharp}}\arrow[r,dash]\arrow[d,hook]\&
Z_{\rho,\sigma}\arrow[r,"\delta_Z"]\arrow[d,hook]\&
\gamma_{\rho,\sigma}\arrow[r]\arrow[d,Rightarrow]\&1\\
1\arrow[r,dash]\&
\mathcal{G}^{\unsure{\natural|\sharp}}\arrow[r,dash]\arrow[d,hook]\&
G_{\rho,\sigma}^{\unsure{\natural|\sharp}}\arrow[r,"\delta_G"]\arrow[d,hook]\&
\gamma_{\rho,\sigma}\arrow[r]\arrow[d,Rightarrow]\&1.\\
1\arrow[r,dash]\&\mathcal{G}^{+}\arrow[r,dash]\&
G_{\rho,\sigma}\arrow[r]\&\gamma_{\rho,\sigma}\arrow[r]\&1
\end{tikzcd}
}
On considère les composés
\diagram{
\begin{tikzcd}[ampersand replacement=\&]
\&\&\mu\\
(33)\quad G_{\rho,\sigma}\arrow[r,"\delta_G"]\&
\gamma_{\rho,\sigma}
 \arrow[ur,"\varepsilon_{2\gamma}"]
 \arrow[r,"\chi_\gamma"]
 \arrow[dr,"\varepsilon_{3\gamma}"']\&
\widehat{\mathbb{Z}}^{*}\\
\&\&\mu
\end{tikzcd}
}
notés
\[
\deleted{\unsure{\varepsilon_G=\varepsilon_\gamma\delta_G}}
\]
\[
\chi_G=\chi_\gamma\circ\delta_G,\qquad
\varepsilon_{2G}=\deleted{\unsure{\chi}}\inserted{\varepsilon}_{2\gamma}\delta_G,\qquad
\varepsilon_{3G}=\delta_G\varepsilon_{3\gamma}.
\]

On a maintenant
\[
(34)\qquad
\left\{
\begin{array}{l}
G'_{\rho,\sigma},G''_{\rho,\sigma},G'''_{\rho,\sigma},G^\circ_{\rho,\sigma}
\quad\text{comme dans }(24)\\
G^+_{\rho,\sigma}=\operatorname{Ker}\chi_G,\quad
G^{'+}_{\rho,\sigma},G^{''+}_{\rho,\sigma},G^{'''+}_{\rho,\sigma},G^{\circ+}_{\rho,\sigma}\\
(\gamma^+_{\rho,\sigma}=\operatorname{Ker}\chi_\gamma\quad
\text{\unsure{noté}}\ \ldots\unsure{\ldots}\quad)
\end{array}
\right.
\]
\margin{idem \unsure{avec} $G_{\rho,\sigma}^{'\unsure{\natural|\sharp}}\ldots$}
Si
\[
(35)\qquad
\tau_G\in G_{\rho,\sigma}^{\unsure{\natural|\sharp}}
\inserted{\subset G_{\rho,\sigma}},\quad
\tau_\gamma\in\gamma_{\rho,\sigma}
\]
\mspage{659}

\scanpage{661}
\mspage{633}
désignent les images de $\tau_Z$ (26) \unsure{ds|dans}
$G_{\rho,\sigma},\gamma_{\rho,\sigma}$,
on peut faire l'identification
\[
(36)\qquad
\left\{
\begin{array}{l}
\mathcal G\simeq\text{image inverse de }\{1,\tau_\gamma\}
\inserted{\subset\gamma_{\rho,\sigma}}\deleted{\text{dans}}\\
\qquad G_{\rho,\sigma},\quad\text{par }\unsure{G_{\rho,\sigma}}
\longrightarrow\gamma_{\rho,\sigma},\\
\mathcal G^{\unsure{\natural|\sharp}}
=\text{image inv. de }\{1,\tau_\gamma\}\subset\gamma_{\rho,\sigma}
\text{ dans }G_{\rho,\sigma}^{\unsure{\natural|\sharp}},
\end{array}
\right.
\]
et
\[
(37)\qquad
\varepsilon_{\mathcal G}
=\chi_G|_{\mathcal G}
=\varepsilon_{2G}|_{\mathcal G}
=\varepsilon_{3G}|_{\mathcal G}
:\left\{
\begin{array}{l}
\text{triviale sur }\mathcal G^+\\
-1\text{ sur }\tau.
\end{array}
\right.
\]

Considérons l'application canonique
\[
(38)\qquad
\diagram{
\begin{tikzcd}[ampersand replacement=\&, column sep=small,row sep=large]
G_{\rho,\sigma}^{\unsure{\natural|\sharp}}
\arrow[r,hook] \&
G_{\rho,\sigma}^{\deleted{\unsure{\natural|\sharp}}}
\arrow[r] \&
\operatorname{Aut}(\mathcal G^+)\times\widehat{\mathbb Z}^{*}\times\mu\times\mu\\
\&\uline{u}\arrow[r]\&
\bigl(\operatorname{int}(\uline{u})|_{\mathcal G^+},
\chi_G(u),\varepsilon_{2G}(u),\varepsilon_{3G}(u)\bigr)
\end{tikzcd}
}
\]

\uline{Proposition}\quad C'est un hom. de groupes:
\deleted{\unsure{$G_{\rho,\sigma}$}}\deleted{dont}
l'image \inserted{de $G_{\rho,\sigma}^{\unsure{\natural|\sharp}}$}
est formée des
\deleted{\unsure{couples}}\inserted{\unsure{quadruples|quadrupliens}}
\[
(u,\mu,\varepsilon_2,\varepsilon_3)\in
\operatorname{Aut}(\mathcal G^+)\times\widehat{\mathbb Z}^{*}\times\mu\times\mu
\]
satisfaisant les conditions
\[
(39)\qquad
\left\{
\begin{array}{ll}
\text{a)}&
u(\rho)\text{ est }\mathcal G_\tau^{\unsure{\natural|\sharp}}
\text{-conj. à }\varepsilon_3[\tau]\rho\\
\text{b)}&
u(\sigma)\quad\rule{4.2cm}{0.35pt}\quad
\unsure{\varepsilon_2|\varepsilon_\tau}[\tau]\sigma\\
\text{c)}&
u(\varepsilon_0)\text{ est }\mathcal G^{\unsure{\natural|\sharp}}
\text{-conjugué à }\varepsilon_0^\mu\\
\text{d)}&
u\text{ stabilise }\mathcal G^{\unsure{\natural|\sharp}}
\quad(\text{automatique vu a), b) si }\varepsilon_2=\varepsilon_3)
\end{array}
\right.
\]
Son noyau est \deleted{isomorphe au}
\deleted{\inserted{formé des éléments}}
\[
(40)\quad
\deleted{\unsure{f^{-1}\uline{u}}}
\left\{
f^{-1}\uline{u}\in G_{\rho,\sigma}^{\unsure{\natural|\sharp}}
\ \middle|\
\begin{array}{l}
\deleted{\unsure{f}}\uline{u}=(u,1,1,1)\in Z_{\rho,\sigma}^{\circ+}\\
f\in\mathcal G^{\unsure{\natural|\sharp}}
\end{array}
\quad\left\{\vphantom{\begin{array}{l}u\\u\end{array}}
u=\operatorname{int}(f).
\right.
\right\}.
\]

\deleted{
\uline{Corollaire}\quad Soit
\[
\unsure{E|\mathcal E}\qquad
(41)\quad
\overset{\inserted{\unsure{\operatorname{Ker}}}}{\unsure{Z_0}}
=\{f\in\mathcal G^{\unsure{\natural|\sharp}}\mid
f(\rho)\text{ est }\mathcal G_\tau^{\unsure{\natural|\sharp}}
\text{-conjugué à }\rho\unsure{\}}.
\]
\[
Z_0^{\unsure{\natural|\sharp}}
=\operatorname{Centr}_{\unsure{\mathcal G^{\natural}|\mathcal G}}
(\varepsilon_0\unsure{)}.
\]
}

\uline{Corollaire}\quad Soient
\[
(41)\qquad
Z_0=\operatorname{Centr}_{\mathcal G^+}(L_0),\qquad
Z_0^{\unsure{\natural|\sharp}}
=Z_0\cap\mathcal G^{\unsure{\natural|\sharp}}
\]
\[
\inserted{
\left(=\{f\in\mathcal G^{\unsure{\natural|\sharp}}\mid
(\operatorname{int}(f),1,1,1)\in
\unsure{Z_{\rho,\sigma}^{\circ+}|Z_{\rho,\sigma}}\unsure{\}}\unsure{)}\right.
}
\]
de sorte que $L_0^{\unsure{\natural|\sharp}}$ s'identifie à un sous-groupe
\mspage{660}

\scanpage{662}
central de $Z_0^{\unsure{\natural|\sharp}}$.
Considérons l'application
\[
(42)\qquad
\diagram{
\begin{tikzcd}[ampersand replacement=\&, column sep=small,row sep=large]
Z_0^{\unsure{\natural|\sharp}}\arrow[rr]\&\&
S_0G_{\rho,\sigma}^{\unsure{\natural|\sharp}}
=Z_{\rho,\sigma}\cdot\mathcal G^{\unsure{\natural|\sharp}}
\&\&\\
f\arrow[r,mapsto]\&
f^{-1}\arrow[d,phantom,"\in"{rotate=-90}]\&
{\cdot\uline{u}(f)}\arrow[d,phantom,"\in"{rotate=-90}]\&
\inserted{=\uline{u}(f)\cdot f^{-1}}\&
{\text{où }\uline{u}(f)=(\operatorname{int}(f),1,1,1)
\in Z_{\rho,\sigma}}\arrow[d,phantom,"\subset"{rotate=-90}]\\
\&\mathcal G^{\unsure{\natural|\sharp}}\&
Z_{\rho,\sigma}\&\&G_{\rho,\sigma}^{\unsure{\natural|\sharp}}
\end{tikzcd}
}
\]
Cette application est un hom de groupes,
\deleted{\unsure{dont|déf.}}\inserted{elle}
\deleted{\inserted{image inverse de
$\unsure{Z_0|L_0}\ \deleted{\unsure{L_0}}$}}
\deleted{\unsure{injectif} \unsure{par}
\unsure{proj.|passage} \unsure{au} \unsure{quotient}
$L_0^{\unsure{\natural|\sharp}}$}
\deleted{l'image}
induit sur $L_0^{\unsure{\natural|\sharp}}\subset
Z_0^{\unsure{\natural|\sharp}}$ l'hom canonique
$\inserted{\unsure{i_{L_0,SG}|i_{L_0,S_0G}}}$
(28), et $L_0^{\unsure{\natural|\sharp}}\subset
Z_0^{\unsure{\natural|\sharp}}$ est l'image inverse de
\[
\unsure{i_{L_0,SG}|i_{L_0,S_0G}}
(L_0^{\unsure{\natural|\sharp}})
=\operatorname{Ker}
\bigl(S_0G_{\rho,\sigma}^{\unsure{\natural|\sharp}}
\longrightarrow G_{\rho,\sigma}^{\unsure{\natural|\sharp}}\bigr).
\]
On trouve ainsi un hom injectif
\[
(43)\qquad
\diagram{
\begin{tikzcd}[ampersand replacement=\&]
Z_0^{\unsure{\natural|\sharp}}/L_0^{\unsure{\natural|\sharp}}
\arrow[r,hook]\&G_{\rho,\sigma}^{\unsure{\natural|\sharp}}
\end{tikzcd}
}
\]
dont l'image est le noyau de (38), i.e on a une suite exacte canonique
\[
(44)\qquad
\diagram{
\begin{tikzcd}[ampersand replacement=\&, column sep=small,row sep=large]
1\arrow[r]\&
L_0^{\unsure{\natural|\sharp}}\arrow[r]\&
Z_0^{\unsure{\natural|\sharp}}\arrow[r]\&
G_{\rho,\sigma}^{\unsure{\natural|\sharp}}\arrow[r]\&
\widetilde G_{\rho,\sigma}^{\unsure{\natural|\sharp}}
\arrow[r]\arrow[d,phantom,"\subset"{rotate=-90}]\&1\\
\&\&\&\&
{\operatorname{Aut}(\mathcal G^+)\times\widehat{\mathbb Z}^{*}\times\mu\times\mu}
\&
\end{tikzcd}
}
\]
où $\widetilde G_{\rho,\sigma}^{\unsure{\natural|\sharp}}$
est formé des $(u,\mu,\varepsilon_2,\varepsilon_3)$ satisfaisant
les conditions (39).

\uline{Remarque.}\quad Cette \unsure{projection|présentation} fait apparaître
dans quelle mesure la construction faite ici de
$G_{\rho,\sigma}^{\unsure{\natural|\sharp}}$
est plus fine que celle \unsure{où}
on le \unsure{définirait|définissait} simplement comme un
groupe d'autom. de $\mathcal G^+$
(on $\overline{\mathrm m}$ \unsure{a|en} besoin
des $\mu,\varepsilon,\varepsilon'$
pr \unsure{desserrer|donner} le \unsure{moule|modèle}\ldots).
\mspage{661}
